\section{Introduction}
\paragraph{Motivation and Context.} 
The phenomenon of buckling of thin-walled structures (shell) describes the transition from 
an initial unstable configuration to a possibly large-magnitude
 displacement under small increase in loading. This so-called buckling response 
 often causes either a significant reduction in stiffness or a complete structural failure.
Relevant examples of shells that are prone to buckling are lightweight air- or 
spacecrafts, submarines and automotives \cite{WAGNER2020106451}. 
Notoriously, in a major series of laboratory experiments by NASA in the 1970s 
\cite{NASA_SP_8007_2020}, the knockdown factor, 
the ratio between observed and theoretical critical buckling load, was shown 
to be as low as twenty per cent. 
Despite the advances in computational resources, mathematical models and 
manufacturing technology, this report remains the most prominent source of 
recommendations for the design of buckling-critical shells. 

%And thus, one motivation is to improve such manufacturing specifications 
%by performing forward and inverse uncertainty quantification for the critical buckling load.
%sharpening 
%the bound on the critical bucklo through a general treatment of domain uncertainty
%propagated to the critical buckling load.
%and extending the quantification of sensitivity to 
%different geometries and geometric imperfections. 
It is well recognized that geometric imperfections are the primary 
reason for the large discrepancy between the observed and 
theoretical critical buckling load \cite{WAGNER2020106451}. 
Naturally, efforts have been directed to the theoretical and 
numerical analysis of geometric imperfection sensitivity of the critical buckling model 
\cite{HUTCHINSON2018157,Harutyunyan2023,HILBURGER2004369, SCHENK20073424}. 
In this article, we apply a general uncertainty quantification analysis of 
the generalized eigenvalue problem that results from the linear buckling model of general shell structures. 
Moreover, we show Gevrey-regularity of the corresponding 
critical buckling load under the assumption that the random domain perturbations are 
Gevrey-regular and uniformly diffeomorphic. The accompanied 
constants are central to inform higher-order convergence methods 
such as quasi Monte-Carlo or 
sparse grid approximation methods \cite{kuo2016applicationquasimontecarlomethods,BieriDiss, Haji_Ali_2016}. 
Consequently, this result enables standard 
applications of methods in forward uncertainty quantification, 
inverse problems, and topological and structural as well as 
experimental design optimization \cite{kaarnioja2026quasimontecarlobayesiandesign,LatifiRostami2019}.


\paragraph{Problem statement.} 
The eigenvalue problem associated to the buckling problem is obtained from considering a 
linearization of the second variation of the elastic energy functional. 
With an increase in load, initially stable configurations may 
change sign in the second variation of energy and thus become unstable. 
Subsequently, a structure in an unstable configuration can suddenly 
transition to a large-magnitude deformation. In this linearization 
the pre-stress components involved in the non-linear term of the second 
variation of elastic energy depend linearly on the traction 
load and can hence be factored out as a dimensionless constant. The 
resulting equations resemble a general eigenvalue problem where 
the lowest eigenvalue corresponds to the dimensionless critical buckling load. 
Here, the considered model is based on a geometric non-linear 
deformation energy function and was analyzed, for example, in 
\cite{DaugeandSuriNumericalapproximation, SZABO199943}.

The parametrized general eigenvalue problem reads as:
\emph{Find the smallest value $\lambda(\bsy) $ and $\bsu(\bsy)$ such that}
\begin{align}\label{eq:bucklingPreview}
    a_{\bsy}(\bsu(\bsy), \bsv) = \lambda(\bsy) b_{\bsy}(\bsu, \bsv) , \quad  \forall \bsv \in V,
\end{align}
where $a_\bsy$ is the variational bilinear form for linear elasticity, $b_\bsy$ is the pre-stress bilinear form and 
$\bsy \in \calU$ is a parameter for the random domain mapping that models domain uncertainty in the shell. 
Perturbations in the domain are modeled with the framework of random domain mappings \cite{Harbrecht2023DomainUncertainty} 
and enter the equation through a pull-back to the nominal domain.
Since parametrized eigenvalues are known to admit non-smooth dependencies, 
we need to restrict our analysis on eigenvalue clusters. 
Eigenvalue cluster are subsets of eigenvalues that remain isolated; that is, 
the eigenvalues of the corresponding eigenmodes do 
not cross with the eigenvalues of eigenmodes outside of the cluster for the whole parameter domain. 
We are concerned with the question of whether Gevrey-regularity of the 
random domain is preserved for the eigenvalue cluster, 
in the sense that its spectral operator is Gevrey-regular. 
Secondly, we want to quantify how large the eigenvalue cluster is 
and if it allows for an approximation of the critical buckling load.
 
\paragraph{State of the art and the Research gap.}
Parametrized eigenvalue problems are well studied, see for example 
\cite{grubišić2021stochasticcollocationmethodcomputing,chernov2024analyticgevreyclassregularity} 
and the references therein. However, most existing results are only valid for compact operators for which the whole spectrum 
is discrete and where the eigenvalue of interest is isolated. The buckling problem, however, admits a few challenges: 
With the Riesz representation theorem on \eqref{eq:bucklingPreview}, we obtain the linear operator general 
eigenvalue problem
\begin{align*}
    \calA_\bfy \bfu &= \lambda(\bfy)\calM_{\bfy} \bfu,
\end{align*}
where $\calA_\bsy$ is uniformly invertible and $\calM_\bsy$ is possibly indefinite.  
Moreover, since $a_\bsy$ and $b_\bsy$ are a vector-valued first order bilinear forms, 
the spectral operator $\calT_\bsy = \calA_\bsy^{-1} \calM_\bsy$ is non-compact. It is known by 
\cite{DaugeandSuriNumericalapproximation} that the non-compact spectral operator from \eqref{eq:bucklingPreview} 
is prone to spectral pollution; whole intervals can become eigenvalues and thus numerical approximations 
seem to not converge in a certain range of eigenvalues. In \cite{DaugeandSuriNumericalapproximation}
has shown that the smallest eigenvalue can be approximated pollution-free whenever the structure 
is sufficiently thin. 

Previous work on the buckling phenomenon has focused on the propagation of the effect of small idealized perturbations, 
based on analytic extensions of the dimension reduced model solution, on the critical buckling 
load in  \cite{HUTCHINSON2018157}. In related work \cite{Harutyunyan2023}, 
the authors proved that the observed asymptotic 
behavior of critical buckling load changes 
under curvature variations and that the critical load decays 
with a larger rate than the theoretic rate for a perfect geometry 
when the thickness tends to zero. % \cite{dauge:hal-00012083}.
On the numerical side, in \cite{SCHENK20073424} the authors have conducted extensive 
numerical experiments with thin-walled cylinders with geometric imperfections that matches the manufacturing specimens
made available by the initial imperfection data bank in \cite{arbocz1979initial}. It is shown in \cite{HILBURGER2004369} 
that the prediction gap can be as low as 5\% percent when geometric imperfections are included.
To the best of our knowledge, general geometric imperfections have not been considered 
theoretically in the context of buckling.
%Here, we analyze the propagation of random domain perturbations through the associated linear eigenvalue
 %problem using the general uncertainty quantification methodological framework.

\paragraph{Contribution.}
 %We analyze the propagation of random domain perturbations through the associated linear eigenvalue
%problem using the general uncertainty quantification methodological framework. 
Our main result is that the
eigenvalue cluster of the smallest eigenvalue, that is associated to the critical buckling load, 
preserves Gevrey-regularity with an explicitly provided constant.
%when the domain perturbation function is Gevrey-regular 
%and its total variation is bounded 
We introduce several new proof techniques, 
further leveraging the alternating factorial 
technique \cite{chernov2024analyticgevreyclassregularity}. 
We show Gevrey-regularity of the spectral operator by differentiating
its resolvent \cite{kato1995perturbation}. 
%The resolvent method can be used for the 
%computation of the eigenspace and further exploited for the spectral analysis.

Another open question is the asymptotic analysis 
for the critical buckling load, that is, the 
behaviour of the spectrum when the thickness of 
the structure goes to zero. In this limit case, 
we observe in numerical simulation that mixing 
of eigenmodes with high wave numbers with the 
smallest eigenvalue occurs. It remains outside 
of the scope of this work to combine the 
asymptotic analysis with the parametrized 
general eigenvalue problem.






\paragraph{Outline.}
The linear buckling model with domain perturbation is described in Section \ref{sec:LinearModel}. 
In Section \ref{section:AbstractResults}, we prove abstract results on the 
spectrum of $\calT_{\bfy}$ and particularly, prove that an isolated 
eigenvalue cluster of $\calT_\bfy$ preserves Gevrey-regularity. %provided that  $\bfy \mapsto \calT_\bfy$ is Gevrey-regular. 
In Section \ref{sec:regularityOfTy}, we show that $\calT_\bfy$ is 
Gevrey-regular if the random domain mapping is Gevrey-regular and uniformly 
diffeomorphic. Morevover, we provide the constants of the regularity statement. 
We conclude with numerical experiments in Section 
\ref{sec:numerical-example} and a discussion in Section \ref{sec:Discussion}.




